\documentclass[letterpaper,10pt,english]{article}
\usepackage{%
amsfonts,%
amsmath,%
amssymb,%
amsthm,%
babel,%
bbm,%
%biblatex,%
caption,%
centernot,%
color,%
enumerate,%
%enumitem,%
epsfig,%
epstopdf,%
etex,%
fancybox,%
framed,%
fullpage,%
%geometry,%
graphicx,%
hyperref,%
latexsym,%
mathptmx,%
mathtools,%
multicol,%
paralist,%
pgf,%
pgfplots,%
pgfplotstable,%
pgfpages,%
proof,%
psfrag,%
%subfigure,%
tcolorbox,%
tfrupee,%
tikz,%
times,%
%ulem,%
url,%
xcolor,%
mathpazo,%
}

\definecolor{shadecolor}{gray}{.95}%{rgb}{1,0,0}
\usepackage[margin=1in,top=0.75in]{geometry}
\usepackage[mathscr]{eucal}
\usepgflibrary{shapes}
\usepgfplotslibrary{fillbetween}
\usetikzlibrary{%
arrows,%
backgrounds,%
chains,%
decorations.pathmorphing,% /pgf/decoration/random steps | erste Graphik
decorations.text,% 
matrix,%
positioning,% wg. " of "
fit,%
patterns,%
petri,%
plotmarks,%
scopes,%
shadows,%
shapes.misc,% wg. rounded rectangle
shapes.arrows,%
shapes.callouts,%
shapes%
}

%\pgfplotsset{compat=newest} %<------ Here
\pgfplotsset{compat=1.11} %<------ Or use this one

\theoremstyle{plain}
\newtheorem{thm}{Theorem}[section]
\newtheorem{lem}[thm]{Lemma}
\newtheorem{prop}[thm]{Proposition}
\newtheorem{cor}[thm]{Corollary}
\newtheorem{clm}[thm]{Claim}

\theoremstyle{definition}
\newtheorem{defn}[thm]{Definition}
\newtheorem{conj}[thm]{Conjecture}
\newtheorem{exmp}[thm]{Example}
\newtheorem{axiom}[thm]{Axiom}
\newtheorem{assum}[thm]{Assumption}
\newtheorem{exerc}[thm]{Exercise}
\newtheorem*{defin}{Definition}

\theoremstyle{remark}
\newtheorem{rem}{Remark}
\newtheorem{note}{Note}

\newcommand{\iid}{\emph{i.i.d.}\ }
\newcommand{\Cov}{\operatorname{Cov}}
%\newcommand{\det}{\operatorname{det}}
%\newcommand{\Real}{\mathbb{R}}
\newcommand{\tr}{\operatorname{tr}}
\newcommand{\Var}{\operatorname{Var}}
\newcommand{\cov}{\operatorname{cov}}

%\renewcommand{\proof}[1]{\begin{proof}#1\end{proof}}
\newcommand{\EQ}[1]{\begin{equation*}#1\end{equation*}}
\newcommand{\EQN}[1]{\begin{equation}#1\end{equation}}
\newcommand{\eq}[1]{\begin{align*}#1\end{align*}}
\newcommand{\meq}[2]{\begin{xalignat*}{#1}#2\end{xalignat*}}
\newcommand{\norm}[1]{\left\lVert#1\right\rVert}
\newcommand{\inner}[1]{\left\langle #1\right\rangle}
\newcommand{\abs}[1]{\left\lvert#1\right\rvert}
\newcommand{\expect}[1]{\mathbb{E}\left[{#1}\right]}
\newcommand{\prob}[1]{\mathbb{P}\left[{#1}\right]}
\newcommand{\given}{\; \big\vert \;} 
\newcommand{\set}[1]{\left\{#1\right\}} 
\newcommand{\indicator}[1]{1_{\set{#1}}} 
\newcommand{\Ind}[1]{\mathbbm{1}_{#1}} 
\newcommand{\SetIn}[1]{\mathbbm{1}_{\set{#1}}} 
\newcommand{\red}[1]{\textcolor{red}{#1}} 
\newcommand{\floor}[1]{\left\lfloor#1\right\rfloor}
\newcommand{\ceil}[1]{\left\lceil#1\right\rceil}


\newcommand{\C}{\mathbb{C}}
\newcommand{\D}{\mathbb{D}}
\newcommand{\E}{\mathbb{E}}
\newcommand{\F}{\mathbb{F}}
\newcommand{\N}{\mathbb{N}}
\renewcommand{\P}{\mathbb{P}}
\newcommand{\Q}{\mathbb{Q}}
\newcommand{\R}{\mathbb{R}}
\newcommand{\Z}{\mathbb{Z}}

\newcommand{\bA}{\mathbf{A}}
\newcommand{\bI}{\mathbf{I}}
\newcommand{\bU}{\mathbf{1}}
\newcommand{\bx}{\mathbf{x}}
\newcommand{\bX}{\mathbf{X}}
\newcommand{\bZ}{\mathbf{Z}}


\newcommand{\cA}{\mathcal{A}}
\newcommand{\cB}{\mathcal{B}}
\newcommand{\cC}{\mathcal{C}}
\newcommand{\cF}{\mathcal{F}}
\newcommand{\cG}{\mathcal{G}}
\newcommand{\cM}{\mathcal{M}}
\newcommand{\cN}{\mathcal{N}}
\newcommand{\cP}{\mathcal{P}}
\newcommand{\cT}{\mathcal{T}}
\newcommand{\cX}{\mathcal{X}}
\newcommand{\cY}{\mathcal{Y}}

\newcommand{\sA}{\mathscr{A}}
\newcommand{\sB}{\mathscr{B}}
\newcommand{\sC}{\mathscr{C}}
\newcommand{\sE}{\mathscr{E}}
\newcommand{\sF}{\mathscr{F}}
\newcommand{\sG}{\mathscr{G}}
\newcommand{\sH}{\mathscr{H}}
\newcommand{\sL}{\mathscr{L}}
\newcommand{\sS}{\mathscr{S}}
\newcommand{\sT}{\mathscr{T}}
\newcommand{\sX}{\mathscr{X}}
\newcommand{\sY}{\mathscr{Y}}
\newcommand{\sZ}{\mathscr{Z}}

\newcommand{\tS}{\tilde{S}}
% Debug
\newcommand{\todo}[1]{\begin{color}{blue}{{\bf~[TODO:~#1]}}\end{color}}

% a few handy macros

\renewcommand{\le}{\leqslant}
\renewcommand{\ge}{\geqslant}
\newcommand\matlab{{\sc matlab}}
\newcommand{\goto}{\rightarrow}
\newcommand{\bigo}{{\mathcal O}}
%\newcommand{\half}{\frac{1}{2}}
%\newcommand\implies{\quad\Longrightarrow\quad}
\newcommand\reals{{{\rm l} \kern -.15em {\rm R} }}
\newcommand\complex{{\raisebox{.043ex}{\rule{0.07em}{1.56ex}} \hskip -.35em {\rm C}}}


% macros for matrices/vectors:

% matrix environment for vectors or matrices where elements are centered
\newenvironment{mat}{\left[\begin{array}{ccccccccccccccc}}{\end{array}\right]}
\newcommand\bcm{\begin{mat}}
\newcommand\ecm{\end{mat}}

% matrix environment for vectors or matrices where elements are right justifvied
\newenvironment{rmat}{\left[\begin{array}{rrrrrrrrrrrrr}}{\end{array}\right]}
\newcommand\brm{\begin{rmat}}
\newcommand\erm{\end{rmat}}

% for left brace and a set of choices
%\newenvironment{choices}{\left\{ \begin{array}{ll}}{\end{array}\right.}
\newcommand\when{&\text{if~}}
\newcommand\otherwise{&\text{otherwise}}
% sample usage:
%\delta_{ij} = \begin{choices} 1 \when i=j, \\ 0 \otherwise \end{choices}


% for labeling and referencing equations:
\newcommand{\eql}{\begin{equation}\label}
\newcommand{\eqn}[1]{(\ref{#1})}
% can then do
%\eql{eqnlabel}
%...
%\end{equation}
% and refer to it as equation \eqn{eqnlabel}.
% some useful macros for finite difference methods:
\newcommand\unp{U^{n+1}}
\newcommand\unm{U^{n-1}}
% for chemical reactions:
\newcommand{\react}[1]{\stackrel{K_{#1}}{\rightarrow}}
\newcommand{\reactb}[2]{\stackrel{K_{#1}}{~\stackrel{\rightleftharpoons}
 {\scriptstyle K_{#2}}}~}
\makeatletter
\def\th@plain{%
\thm@notefont{}% same as heading font
\itshape % body font
}
\def\th@definition{%
\thm@notefont{}% same as heading font
\normalfont % body font
}
\makeatother
\date{}



\title{ Lecture-19: Discrete Time Markov Chains}
\author{}

\begin{document}
\maketitle

\section{Markov processes}
We have seen that \iid sequences are easiest discrete time random processes. 
However, they don't capture correlation well. 
\begin{defn}
A stochastic process $X:\Omega\to\sX^T$ with state space $\sX$ and ordered index set $T$ is \textbf{Markov} if conditioned on the present state $X_t$, future $\sigma(X_u, u > t)$ is independent of the past $\sigma(X_s, s< t)$. 
We denote the history of the process until time $t$ as $\sF_t \triangleq \sigma(X_s, s \le t)$. 
That is, for any Borel measurable set $B \in \sB(\sX)$ %any ordered index set $T$ containing 
any two indices $u > t$, we have  %such that $t < u$, % \in T$, 
\EQ{
P(\set{X_u \in B}\mid \sF_t) = P(\set{X_{u} \in B}\mid \sigma(X_t)).
}
%The range of the process is called the \textbf{state space}. 
\end{defn}
\begin{rem}
%For a stationary Markov process, we would have 
We next re-write the Markov property more explicitly for the process $X:\Omega\to\R^\R$. 
For all $x, y \in \sX$, finite set $S \subseteq \R$ such that $\max S < t < u$, and $H_S(x_S) \triangleq \cap_{s \in S}\set{X_s \le x_s} \in \sF_t$, we have 
\EQ{
P(\set{X_u \le y}\mid H_S(x_S) \cap \set{X_t \le x}) = P(\set{X_{u} \le y}\mid \set{X_t \le x}).
}
\end{rem}
%\begin{rem}
%When the state space $\sX$ is countable, we can write $H_S = \cap_{s \in S}\set{X_s = x_s}$ and the Markov property can be written as 
%\EQ{
%P(\set{X_u = y}\mid H_S \cap \set{X_t = x}) = P(\set{X_{u} = x_u}\mid \set{X_t = x}).
%}
%\end{rem}
%\begin{rem}
%In addition, consider the case when the index set is countable, i.e. $T = \Z_+$. 
%Then, for any present instant $n \in \Z_+$, we can take past as $\sF_{n-1}$, and the future state as $X_{n+1}$.
%Then, the Markov property for any states $x,y\in\sX$ can be written as 
%\EQ{
%P(\set{X_{n+1} = y}\mid H_{n-1} \cap \set{X_n = x}) = P(\set{X_{n+1} = y}\mid \set{X_n = x}).
%}
%\end{rem}
%We will study this process in detail in coming lectures. 
%\begin{shaded}
%%\begin{exmp}
%%A random walk $S:\Omega\to\sX^{\Z_+}$ with \iid step-size sequence $X:\Omega\to\sX^\N$, is a homogeneous Markov sequence. 
%%For any $n \in \Z_+$ and $x,y, s_1, \dots, s_{n-1} \in \sX$, we can write the conditional probability 
%%\EQ{
%%P(\set{S_{n+1} = y}\mid\set{S_n = x, S_{n-1} = s_{n-1}, \dots, S_1 = s_1}) 
%%= P(\set{S_{n+1} - S_n = y-x})
%%=  P(\set{S_{n+1} = y}\mid\set{S_n = x}).
%%}
%%\end{exmp}
%\begin{lem} 
%The random $S:\Omega\to\Z_+^{\Z_+}$ with \iid step-size sequence $X:\Omega\to\sX^\N$, is homogeneously Markov. 
%\end{lem}
%\begin{proof} 
%Since the process has stationary and independent increments, we have 
%\EQ{
%P(\set{S_{n+m} = k} \mid\set{S_1 = k_1, S_2  = k_2, \dots, S_n = k_n}) 
%= P(\set{S_{n+m}  - S_n = k - k_n}) 
%= P(\set{S_{n+m} = k}\mid\set{S_n = k_n}). 
%}
%\end{proof}
%\end{shaded}


\subsection{Discrete time Markov chains}
\begin{defn}
For a state space $\sX \subseteq \R$ and the random sequence $X: \Omega \to \sX^{\Z_+}$, 
we define the history until time $n \in \Z_+$ as $\sF_n \triangleq \sigma(X_1, \dots, X_n)$. 
\end{defn} 
\begin{rem} 
Recall that the event space $\sF_n$ is generated by the historical events $A_X(x) \triangleq \cap_{i=0}^n\set{X_i \le x_i}$ where $x \in \R^{n+1}$.
\end{rem}
\begin{rem} 
When the state space $\sX$ is countable, the event space $\sF_n$ is generated by the historical events $H_n(x) \triangleq \cap_{i=0}^n\set{X_i = x_i}$, where $x \in \sX^{n+1}$. 
That is, $\sF_n = \sigma(H_n(x):x \in \sX^{n+1})$
\end{rem}
%%If $X_n \in S$ for all $n \in \Z_+$, then $E$ is called the \textbf{state space} of process $X$. 
%%We consider a countable state space, and we say that the process $X$ is in state $i$ at time $n$ if $X_n = i \in S$.   
\begin{defn} %[DTMC]
For a countable set $\sX$, a discrete-valued random sequence $X: \Omega \to \sX^{\Z_+}$ is called a \textbf{discrete time Markov chain (DTMC)} 
if for all positive integers $n \in \Z_+$, all states $x,y \in \sX$, 
and any historical event $H_{n-1} = \cap_{m=0}^{n-1}\set{X_m = x_m} \in \sF_n$ for $(x_0, \dots, x_{n-1}) \in \sX^{n}$, 
the process $X$ satisfies the Markov property 
\EQ{
%P(\set{X_{n+1} = y}|\sigma(X_0, \dots, X_n)) = P(\set{X_{n+1} = y}|\sigma(X_n)). 
P(\set{X_{n+1} = y}\given H_{n-1}\cap\set{X_n = x}) = P(\set{X_{n+1} = y}\given \set{X_n = x}). 
}
%where $\sF_n = \sigma(X_0, \dots, X_n)$ %= \sigma(Z_0, \dots, Z_n)$ 
%is the natural filtration. 
%Since the state space $S$ is countable, the probability $P(\set{X_{n+1} = j}|\sigma(X_n))$ can be written as
%\EQ{
%P(\set{X_{n+1} = j}|\sigma(X_n)) = \sum_{i \in S}1_{\set{X_n = i}}P(\set{X_{n+1} = j}|\sigma(X_n)) =  \sum_{i \in S}1_{\set{X_n = i}}P(\set{X_{n+1} = j}|\set{X_n = i}). 
%}
\end{defn}
\begin{rem} 
The above definition is equivalent to $P(\set{X_{n+1} \le x}\given \sF_n) = P(\set{X_{n+1} \le x}\given \sigma(X_n))$ for discrete time discrete state space Markov chain $X$,   
since $\sF_n = \sigma(H_n(x): x \in \sX^n)$ and $\sigma(X_n) = \sigma(\set{X_n = x}, x \in \sX)$. 
\end{rem}
\begin{shaded}
\begin{exmp}[Random Walk] 
A random walk $S:\Omega\to\sX^\N$ with independent step-size sequence $X:\Omega\to\sX^\N$, 
is Markov for a countable state space $\sX$ that is closed under addition. 
%For the countable state space $\sX$, an 
Given a historical event $H_{n-1}(s) \triangleq \cap_{k=1}^{n-1}\set{S_k = s_k}$ and the current state $\set{S_n = s_n}$, 
we can write the conditional probability 
\eq{
P(\set{S_{n+1} = s_{n+1}}\given H_{n-1}(s)\cap\set{S_n = s_n}) 
&= P(\set{X_{n+1} = s_{n+1}-S_n}\given H_{n-1}(s)\cap\set{S_n = s_n})\\
&= P(\set{S_{n+1} = s_{n+1}}\given \set{S_n = s_n}) 
= P\set{X_{n+1} 
= s_{n+1}-s_n} .
}
The equality in the second line follows from the independence of the step-size sequence. 
In particular, from the independence of $X_{n+1}$ from the collection $\sigma(S_0, X_1, \dots, X_n) = \sigma(S_0, S_1, \dots, S_n)$. 
\end{exmp}
\end{shaded}

\subsection{Transition probability matrix}
\begin{defn}
We denote the set of all probability mass functions over a countable state space $\sX$ by 
$
\cM(\sX)\triangleq \set{\nu \in [0,1]^\sX: \sum_{x\in \sX}\nu_x = 1}. 
$
\end{defn}
\begin{defn} 
The \textbf{transition probability matrix} at time $n$ is denoted by $P(n) \in [0,1]^{\sX \times \sX}$, 
such that its $(x,y)$th entry is denoted by $p_{xy}(n) \triangleq P(\set{X_{n+1} = y}\given\set{X_n = x})$, that is the \textbf{transition probability} of a discrete time Markov chain $X$ from a state $x \in \sX$ at time $n$ to state $y \in \sX$ at time $n+1$. 
\end{defn}
\begin{rem}
We observe that each row $P_{x}(n) \triangleq (p_{xy}(n): y \in \sX) \in\cM(\sX)$ is the conditional distribution of $X_{n+1}$ given the event $\set{X_n = x}$. 
\end{rem}
%\begin{defn}[Random walk] 
%Consider the state space $\sX = \R^d$, and the countable index set $T = \Z_+$. 
%Consider the random process $X: \Omega \to \sX^{\Z_+}$ be an be an independent (not necessarily identical) sequence. 
%Let $S_0 = 0$ and $S_n \triangleq \sum_{i=1}^nX_i$ for all $n \in \Z_+$, 
%then the process $S: \Omega \to \sX^{\Z_+}$ is called a \textbf{random walk}. 
%\end{defn} 
%\begin{rem} 
%We can think of $S_n$ as the random location of a particle after $n$ steps, 
%where the particle starts from origin and takes steps of size $X_i$ at the $i$th step. 
%\end{rem}
%\begin{shaded}
%%\begin{exmp}[Random Walk] 
%%Consider the state space $\sX = \R^d$, and the countable index set $T = \Z_+$. 
%%Consider the random process $X = (X_n \in \sX^\Omega: n \in \Z_+)$ be an be an independent (not necessarily identical) sequence. 
%%Let $S_0 = 0$ and $S_n \triangleq \sum_{i=1}^nX_i$, then the process $S = (S_n \in \R^d: n \in \N_0)$ is called a \textbf{random walk}. 
%%We can think of $S_n$ as the random location of a particle after $n$ steps, 
%%where the particle starts from origin and takes steps of size $X_i$ at the $i$th step. 
%%From previous section, we know following properties of random walks. 
%\begin{thm}[Random walk]  
%For a random walk $S:\Omega\to\sX^\N$ with independent step-size sequence $X:\Omega\to\sX^\N$, the following are true. 
%\begin{enumerate}[i\_]
%\item The first two moments are $\E S_n = \sum_{i=1}^n\E X_i$ and $\Var[S_n] = \sum_{i=1}^n\Var[X_i]$. 
%\item Random walk is non-stationary with independent increments. 
%The disjoint increments are stationary if the step-size sequence $X$ is identically distributed. 
%\item Random walk is a Markov sequence for countable state space $\sX$. 
%\end{enumerate}
%\end{thm}
%\proof{
%Results follow from the independence of the step-size sequence $X$. 
%\begin{enumerate}[i\_]
%\item Follows from the linearity of expectation and independence of step sizes. 
%\item Since the mean is time dependent, random walk is non-stationary process. 
%Independence of increments follows from the independence of step sizes. 
%That is, since $\sF_n = \sigma(S_0, S_1, \dots, S_n) = \sigma(S_0, X_1, \dots, X_n)$ and the collection $(X_{n+1}, \dots, X_m)$ is independent of $\sigma(S_0, X_1,\dots, X_n)$ for all $m > n$. 
%Since $S_{m}- S_n = X_{n+1}+ \dots + X_{m} \in \sigma(X_{n+1}, \dots, X_m)$, 
%we have the independent increments. 
%When the step-sizes are also identically distributed, 
%the joint distributions of  $(X_1, \dots, X_{m-n})$ and $(X_{n+1}, \dots, X_m)$ are identical. 
%This implies the stationarity of increments for \iid step-sizes.
% 
%\item 
%%Given the historical event $H_{n-1} \triangleq \cap_{k=1}^{n-1}\set{S_k \le s_k}$ and the current state $\set{S_n \le s_n}$, 
%%we can write the conditional probability 
%%\eq{
%%P(\set{S_{n+1} \le s_{n+1}}\given H_{n-1}\cap\set{S_n \le s_n}) &= P(\set{X_{n+1} \le s_{n+1}-S_n}\given H_{n-1}\cap\set{S_n \le s_n})\\
%%&= P(\set{S_{n+1} \le s_{n+1}}\given \set{S_n \le s_n}). % = P\set{X_{n+1} = s_{n+1}-s_n} .
%%}
%%The equality in the second line follows from the independence of the step-size sequence. 
%%In particular, from the independence of $X_{n+1}$ from the collection $\sigma(S_0, X_1, \dots, X_n) = \sigma(S_0, S_1, \dots, S_n)$. 
%For the countable state space $\sX$, an given the historical event $H_{n-1}(s) \triangleq \cap_{k=1}^{n-1}\set{S_k = s_k}$ and the current state $\set{S_n = s_n}$, 
%we can write the conditional probability 
%\eq{
%P(\set{S_{n+1} = s_{n+1}}\given H_{n-1}(s)\cap\set{S_n = s_n}) 
%&= P(\set{X_{n+1} = s_{n+1}-S_n}\given H_{n-1}(s)\cap\set{S_n = s_n})\\
%&= P(\set{S_{n+1} = s_{n+1}}\given \set{S_n = s_n}) 
%= P\set{X_{n+1} 
%= s_{n+1}-s_n} .
%}
%The equality in the second line follows from the independence of the step-size sequence. 
%In particular, from the independence of $X_{n+1}$ from the collection $\sigma(S_0, X_1, \dots, X_n) = \sigma(S_0, S_1, \dots, S_n)$. 
%\end{enumerate}
%}
%%\end{exmp}
%\end{shaded}

%\begin{shaded}
%\begin{exmp}[Integer random walk] 
%Let $Z = (Z_n \in \Z: n \in \N)$ be an independent (not necessarily identical) Bernoulli sequence. 
%Let $X_0 = 0$ and $X_n \triangleq \sum_{i=1}^nZ_i$, then the process $X = (X_n \in \Z: n \in \Z_+)$ is called an \textbf{integer random walk}. 
%%We can think of $S_n$ as the random location of a particle after $n$ steps, 
%%where the particle starts from origin and takes steps of size $X_i$ at the $i$th step. 
%%From previous section, we know following properties of random walks. 
%\begin{thm} 
%For a random walk $(S_n: n \in \N)$ with \emph{iid} step-size sequence $X$, the following are true. 
%\begin{enumerate}[i\_]
%\item The first two moments are $\E S_n = n\E X_i$ and $\Var[S_n] = n\Var[X_i]$. 
%\item Random walk is non-stationary with independent and stationary increments. 
%\item Random walk is a discrete time Markov chain. 
%\end{enumerate}
%\end{thm}
%\end{exmp}
%\end{shaded}

\begin{defn}
%For all states $x,y \in \sX$, 
A matrix $A \in \R_+^{\sX \times \sX}$ with non-negative entries is called \textbf{sub-stochastic} if the row-sum $\sum_{y \in \sX} a_{xy} \le 1$ for all rows $x \in \sX$. 
If the above property holds with equality for all rows, then it is called a \textbf{stochastic} matrix. 
If matrices $A$ and $A^T$ are both stochastic, then the matrix $A$ is called \textbf{doubly stochastic}. 
\end{defn}
%\begin{shaded}
\begin{rem} 
We make the following observations for the stochastic matrices. 
\begin{compactenum}[i\_] 
\item Every probability transition matrix $P(n)$ is a stochastic matrix. 
\item All the entries of a sub-stochastic matrix lie in $[0,1]$. 
\item Each row $A_x \triangleq (a_{xy}: y \in \sX)$ of the stochastic matrix $A \in \R_+^{\sX \times \sX}$ belongs to $\cM(\sX)$. %is probability mass function over the state space $\sX$. 
\item Every finite stochastic matrix has a right eigenvector with unit eigenvalue. 
This can be observed by taking $\bU^T = \begin{bmatrix}1 & \dots & 1\end{bmatrix}$ to be an all-one vector of length $\abs{\sX}$. 
Then we see that $A \bU = \bU$, since 
\EQ{
(A\bU)_x 
= \sum_{y \in \sX}a_{xy}\bU_y 
= \sum_{y \in \sX} a_{xy} 
= \bU_x, \text{ for each }x \in \sX. 
}
\item Every finite doubly stochastic matrix has a left and right eigenvector with unit eigenvalue. 
This follows from the fact that finite stochastic matrices $A$ and $A^T$ have a common right eigenvector $\bU$. 
It follows that $A$ has a left eigenvector $\bU^T$. 
\item For a probability transition matrix $P(n)$, we have 
$
\sum_{y \in \sX}f(y)p_{xy}(n) = \E[f(X_{n+1})\mid\set{X_n =x}].  
$ 
\end{compactenum}
\end{rem}
%\end{shaded}

\subsection{Homogeneous Markov chains}
In general, not much can be said about Markov chains with index dependent transition probabilities. 
%Hence, 
We consider the simpler case where the transition probabilities $p_{xy}(n)= p_{xy}$ are independent of the index. 
\begin{defn}
A discrete time Markov chain with the probability transition matrix $P(n)$ that is independent of the index, 
is called \textbf{time homogeneous}. %and call the linear operator $P = (p_{xy}: x,y \in \sX)$ the \textbf{transition matrix}. 
%The transition matrix $P$ is stochastic matrix. 
\end{defn}
\begin{shaded}
%When $X$ is a Bernoulli sequence, with $P\set{X_i = 1} = p = 1 - P\set{X_i = -1}$, 
%the one dimensional random walk $S$ is an integer valued random sequence with unit step-size. 
\begin{exmp}[Integer random walk] 
\label{exmp:IRW}
For a one-dimensional integer valued random walk $S: \Omega \to \Z^\N$ with \iid unit step size sequence $X: \Omega \to \set{-1,1}^\N$ such that $P\set{X_1 = 1} = p$, 
the transition operator $P \in [0,1]^{\Z\times\Z}$ is given by the entries 
$
p_{xy} = p\indicator{y=x+1}+(1-p)\indicator{y=x-1} 
$
for all $x,y \in\Z$. 
\end{exmp}
\begin{exmp}[Sequence of experiments]
\label{exmp:MarkovBerExp}
Consider a random sequence of experiment outcomes $X: \Omega\to\set{0,1}^{\Z_+}$, 
such that $P(\set{X_{n+1} = 0}\mid\set{X_n=0})= 1-q$ and $P(\set{X_{n+1} = 1}\mid\set{X_n=1})= 1-p$ for all $n \in \Z_+$. 
Then, we can write the probability transition matrix as  
\EQ{
P = \begin{bmatrix}1-q & q \\ p & 1-p\end{bmatrix}. 
}
\end{exmp}
\end{shaded}
\begin{defn}
Consider a time homogeneous Markov chain $X: \Omega \to \sX^{\Z_+}$ %= (X_n \in \sX: n \in \Z_+)$ 
with countable state space $\sX$ and transition matrix $P$. 
We would respectively denote the conditional probability of events and conditional expectation of random variables, conditioned on the initial state $\set{X_0 = x}$, by 
\meq{2}{
&P_x
(A) \triangleq P(A \given \set{X_0 = x}),&
&\E_x[Y] \triangleq \expect{Y \given \set{X_0 = x}}.
}
\end{defn}
\begin{prop}
Conditioned on the initial state, any finite dimensional distribution of a homogeneous Markov chain is stationary. 
That is, for any finite $n,m \in \Z_+$ and states $x_0, \dots, x_n \in \sX$, we have 
\EQ{
P(\cap_{i=1}^n\set{X_i = x_i}\given\set{X_0 = x_0}) 
= P(\cap_{i=1}^n\set{X_{m+i} = x_i}\given\set{X_m = x_0}) 
= \prod_{i=1}^np_{x_{i-1}x_{i}}.
}
\end{prop}
\begin{proof}
Consider a homogeneous Markov chain $X:\Omega\to\sX^{\Z_+}$ with natural filtration $\sF_\bullet$ such that $\sF_n \triangleq \sigma(X_0, \dots, X_n)$ for each $n\in\N$. 
%We compute the transition probabilities for the path $(x_1, \dots, x_n)$ taken by 
%\begin{compactenum}[(i)]
%\item the sample path $(X_1, \dots, X_n)$ given the event $\set{X_0 = x_0}$ and 
%\item by the sample path $(X_{m+1}, \dots, X_{m+n})$ given the event $\set{X_m = x_0}$. 
%\end{compactenum}
Using the property of conditional probabilities and Markovity of $X$, 
we can write the conditional probability of sample path $(X_1, \dots, X_n)$ given the event $\set{X_0 = x_0}$ as 
\EQ{
P\Big(\cap_{i=1}^n\set{X_i = x_i}\mid\set{X_0=x_0}\Big)
= \prod_{i=1}^nP\Big(\set{X_i = x_i}\mid\cap_{j=0}^{i-1}\set{X_j=x_j}\Big)
= \prod_{i=1}^nP\Big(\set{X_i = x_i}\mid\set{X_{i-1}=x_{i-1}}\Big).
}
Using the property of conditional probabilities and Markovity of $X$, 
we can write the conditional probability of sample path $(X_{m+1}, \dots, X_{m+n})$ given the event $\set{X_m = x_0}$ as 
\EQ{
P\Big(\cap_{i=1}^n\set{X_{m+i} = x_i}\mid\set{X_m=x_0}\Big)
%= \prod_{i=1}^nP\Big(\set{X_{m+i} = x_i}\mid\cap_{j=1}^{i-1}\set{X_{m+j}=x_j}\Big)
= \prod_{i=1}^nP\Big(\set{X_{m+i} = x_i}\mid\set{X_{m+i-1}=x_{i-1}}\Big).
}
From time-homogeneity of transition probabilities of Markov chain $X$, 
it follows that both the transition probabilities are identical and equal to $\prod_{i=1}^np_{x_{i-1},x_i}$. 
%For each $i \in \set{0, \dots, n}$, we can define events $H_i \triangleq \cap_{j=0}^i\set{X_j = x_j}$. 
%We observe that $H_i = \set{X_i = x_i}\cap H_{i-1}$ and $H_i \in \sF_i$ for all $i \in \N$. 
%From the definition of event $H_{n-1}$ and the conditional probability, 
%we can write 
%\EQ{
%P_{x_0}(H_n) = P_{x_0}(\set{X_n = x_n}\cap H_{n-1}) = P(\set{X_n = x_n}\given H_{n-1})P_{x_0}(H_{n-1}).
%}
%Using  the fact that $H_{n-1} = \set{X_{n-1} = x_{n-1}}\cap H_{n-2}$, and the Markovity and homogeneity of the process $X$, we obtain
%\EQ{
%P(\set{X_n = x_n}\given H_{n-1}) 
%= P(\set{X_n = x_n}\given \set{X_{n-1} = x_{n-1}}\cap H_{n-2}) 
%= p_{x_{n-1}x_n}.
%}
%Inductively, we can write the conditional distribution of $H_{n}$ given the event $\set{X_0 = x_0}$ as $P_{x_0}(H_n) = \prod_{i=1}^np_{x_{i-1}x_i}.$
%Similarly, we can write for the sample path $(X_{m+1}, \dots, X_{m+n})$ given $X_m = x_0$, 
%\EQ{
%P(\set{X_{m+1} = x_1, \dots, X_{m+n} = x_n}\given \set{X_m = x_0}) 
%= \prod_{i=1}^n P(\set{X_{m+i} = x_i)}\given \set{X_{m+i-1} = x_{i-1}}) 
%= \prod_{i=1}^np_{x_{i-1}x_i}.  
%}
\end{proof}
\begin{cor}
The $n$-step transition probabilities are stationary for any homogeneous Markov chain.  
That is,  for any states $x_0,x_n \in \sX$ and $n,m \in \N$, we have 
$
P(\set{X_{n+m} = x_n}\mid\set{ X_m = x_0}) 
=  P(\set{X_{n} = x_n}\mid\set{ X_0 = x_0}). 
$
\end{cor}
\begin{proof}
It follows from summing over all possible paths $(X_0, \dots, X_n)$ and $(X_m, \dots, X_{m+n})$. 
In particular, we can partition events $\set{X_n=x_n}$ and $\set{X_{m+n}=x_n}$ in terms of unions over disjoint paths 
\meq{2}{
&\set{X_n = x_n} = \bigcup_{x\in \sX^{n-1}}\Big(\cap_{i=1}^n\set{X_i = x_i}\Big),&
&\set{X_{m+n} = x_n} = \bigcup_{x\in \sX^{n-1}}\Big(\cap_{i=1}^n\set{X_{m+i}=x_i}\Big).
}
%For each $x \in \sX^{n-1}$, we define events $H_{n-1}(x) \triangleq \cap_{i=1}^{n-1}\set{X_i = x_i}$ and $F_{n-1}(x) \triangleq \cap_{i=1}^{n-1}\set{X_{m+i} = x_i}$. 
%We can partition the events $\set{X_n = x_n}$ and $\set{X_{m+n}=x_n}$  in terms of disjoint events $\set{H_{n-1}(x) \cap\set{X = x_n}: x \in \sX^{n-1}}$ and $\set{F_{n-1}(x)\cap\set{X_{m+n}=x_n}: x \in \sX^{n-1}}$ respectively. 
%That is,
%\meq{2}{
%&\set{X_n = x_n} = \cup_{x\in \sX^{n-1}}H_{n-1}(x)\cap\set{X_n=x_n},&
%&\set{X_{m+n} = x_n} = \cup_{x\in \sX^{n-1}}F_{n-1}(x)\cap\set{X_{m+n} = x_n}.
%}
%From the stationarity in joint distribution conditioned on initial state for the homogeneous Markov chain $X$, we have  $P(F_{n-1}(x)\cap\set{X_{m+n} = x_n} \given \set{X_m = x_0}) 
%= P(H_{n-1}(x)\cap\set{X_n= x_n} \given \set{X_0 = x_0})$. 
%Using the law of total probability, we can write the conditional probability 
%\eq{
%&P_{x_0}\set{X_n = x_n} 
%= \sum_{x \in \sX^{n-1}}P_{x_0}(H_{n-1}(x)\cap\set{X_n= x_n}),\\
%&P(\set{X_{m+n} = x_n}\mid \set{X_m = x_0}) 
%= \sum_{x \in \sX^{n-1}}P(F_{n-1}(x)\cap\set{X_{m+n}=x_n} \mid \set{X_m = x_0}). 
%}
The result follows from the countable additivity of conditional probability for disjoint events of taking distinct paths, 
and the fact that probability of taking same path is identical for both sums. 
\end{proof}
\subsection{Transition graph} 
A time homogeneous Markov chain $X: \Omega\to\sX^\N$ with a probability transition matrix $P$, 
is sometimes represented by a directed weighted graph $G = (\sX, E, w)$, where the set of nodes in the graph $G$ is the state space $\sX$, and the set of directed edges is the set of possible one-step transitions indicated by the initial and the final state, as 
\EQ{
E \triangleq \set{[x,y\rangle \in \sX \times \sX : p_{xy} > 0}. 
}
In addition, this graph has a weight $w_e = p_{xy}$ on each edge $e = [x,y\rangle \in E$. 

\begin{shaded} 
\begin{exmp}[Integer random walk] 
The time homogeneous Markov chain in Example~\ref{exmp:IRW} %For an integer random walk $X = (X_n \in \Z: n \in \N)$ with \iid step-size sequence $Z = (Z_n \in \set{-1,1}, n \in \N)$, 
can be represented by an infinite state weighted graph $G = (\Z, E, w)$, where the edge set is 
\EQ{
E = \set{(n, n+1): n \in \Z}\cup\set{(n, n-1): n \in \Z}.
}
We have plotted the sub-graph of the entire transition graph for states $\set{-1,0,1}$ in Figure~\ref{Fig:RandomWalk}. 
\end{exmp}
\end{shaded}
\begin{figure}[h!]
\centering
\scalebox{1}{\input{../../Figures/RandomWalk}}
\caption{Sub-graph of the entire transition graph for an integer random walk with \iid step-sizes in $\set{-1,1}$ with probability $p$ for the positive step.}
\label{Fig:RandomWalk}
\end{figure}

\begin{shaded} 
\begin{exmp}[Sequence of experiments]
The time homogeneous Markov chain in Example~\ref{exmp:MarkovBerExp} can be represented by the following two-state weighted transition graph $G = (\set{0,1}, E, w)$, 
%Consider the sequence of experiments with the set of outcomes $\sX = \set{0,1}$ with the transition matrix 
%\EQ{
%P = \begin{bmatrix}1-q & q \\ p & 1-p\end{bmatrix}. 
%}
%We have plotted the corresponding transition graph for this two-state Markov chain in 
plotted in Figure~\ref{Fig:TwoStateMC}. 
\end{exmp}
\end{shaded}
\begin{figure}[h!]
\centering
\scalebox{1}{\input{../../Figures/TwoStateMarkovChain}}
\caption{Markov chain for the sequence of experiments with two outcomes.}
\label{Fig:TwoStateMC}
\end{figure}

\subsection{Random mapping theorem}
We saw some example of Markov processes where $X_n = X_{n-1} + Z_n$, 
and $Z:\Omega\to\Lambda^\N$ is an \iid sequence, independent of the initial state $X_0$. 
We will show that any discrete time Markov chain is of this form, 
where the sum is replaced by arbitrary functions.  
\begin{thm}[Random mapping theorem] 
For any DTMC $X:\Omega\to\sX^{\Z_+}$, there exists an \iid sequence $Z \in \Lambda^\N$ and  functions $f_n: \sX \times \Lambda \to \sX$ such that $X_{n} = f_n(X_{n-1}, Z_{n})$ for all $n \in \N$. 
\end{thm}
\begin{rem}
A \textbf{random mapping representation} of a transition matrix $P(n)$ on state space $\sX$ is a function $f_n: \sX \times \Lambda \to \sX$, along with a random variable $Z_n:\Omega\to\Lambda$, satisfying for all $x,y \in \sX$, 
\EQ{
P\set{f_n(x,Z_n) = y} = p_{xy}(n).
}
\end{rem}
\begin{proof}
It suffices to show that every transition matrix $P(n)$ has a random mapping representation. 
Then, for the mapping $f_n$ and the \iid sequence $Z:\Omega\to\Lambda^\N$, we would have $X_n = f_n(X_{n-1}, Z_n)$ for all $n \in \N$. 

Let $\Lambda \triangleq [0,1]$, and we choose the \iid uniform sequence $Z:\Omega\to\Lambda^\N$. %, uniformly at random from this interval. 
Since $\sX$ is countable, it can be ordered. 
We let $\sX = \N$ without any loss of generality. 
We set $F_{xy}(n) \triangleq \sum_{w \le y}p_{xw}(n)$ and define function $f_n:\sX\times\Lambda\to\sX$ for all pairs $(x,z)\in\sX\times\Lambda$ by 
\EQ{
f_n(x, z)  \triangleq \sum_{y \in \N}y\SetIn{ F_{x, y-1}(n)< z \le F_{x,y}(n)} = \inf\set{y \in \sX: z \le F_{x,y}(n)}.
}
Since $f_n(x,Z_n)$ is a discrete random variable taking value $y \in \sX$, iff the uniform random variable $Z_n$ lies in the interval $(F_{x, y-1}(n), F_{x,y}(n)]$. 
That is, the event $\set{f_n(x,Z_n) = y} = \set{Z_n \in (F_{x, y-1}(n), F_{x,y}(n)]}$ for all $y \in \sX$. 
It follows that 
\EQ{
P\set{f_n(x,Z_n) = y} = P\set{F_{x, y-1}(n)< Z_n \le F_{x,y}(n)} =  F_{x,y}(n) - F_{x,y-1}(n) =p_{xy}(n).
} 
\end{proof}
\end{document}
